The Event class represents an instance of an event happening. Each Event
has an EventType that describes the event and pointers to the Process and
Thread that the event occurred on.

The Event class is an abstract class that is never instantiated. Instead,
ProcControlAPI will instantiate children of the Event class, each of which add
information specific to the EventType. For example, an Event representing a
thread creation will have an EventType of ThreadCreate and can be cast into an
EventNewThread for specific information about the new thread.

An event that occurs on a running thread may cause the process, thread, or
neither to stop running until the event has been handled. The specifics of
what is stopped can change between different event types and operating systems.
Each Event describes whether it stopped the associated process or thread with
a SyncType field. The values of this field can be async (the event stopped
neither the process nor thread), sync\_thread (the event stopped its thread),
or sync\_process (the event stopped all threads in the process). A callback
function can choose how to resume or stop a process or thread using its return
value.

\begin{apient}
  enum SyncType { unset, async, sync_thread, sync_process }
\end{apient}

\apidesc{
  The SyncType type is used to describe how a process or thread is stopped by an
  Event. See the above explanation for more details.
}

\begin{apient}
  Thread::const_ptr getThread() const
\end{apient}

\apidesc{
  This function returns a const pointer to the Thread object that represents the
  thread this event occurred on.
}

\begin{apient}
  Process::const_ptr getProcess() const
\end{apient}

\apidesc{
  This function returns a const pointer to the Process object that
  represents the
  process this event occurred on.
}

\begin{apient}
  EventType getEventType() const
\end{apient}

\apidesc{
  This function returns the EventType associated with this Event.
}

\begin{apient}
  SyncType getSyncType() const
\end{apient}

\apidesc{
  This function returns the SyncType associated with this Event.
}

\begin{apient}
  std::string name() const
\end{apient}

\apidesc{
  This function returns a human readable name for this Event.
}

\begin{apient}
  EventTerminate::ptr getEventTerminate()
  EventTerminate::const_ptr getEventTerminate() const
\end{apient}

\apidesc{
}

\begin{apient}
  EventExit::ptr getEventExit()
  EventExit::const_ptr getEventExit() const
\end{apient}

\apidesc{
}

\begin{apient}
  EventCrash::ptr getEventCrash()
  EventCrash::const_ptr getEventCrash() const
\end{apient}

\apidesc{
}

\begin{apient}
  EventForceTerminate::ptr getEventForceTerminate()
  EventForceTerminate::const_ptr getEventForceTerminate() const
\end{apient}

\apidesc{
}

\begin{apient}
  EventExec::ptr getEventExec()
  EventExec::const_ptr getEventExec() const
\end{apient}

\apidesc{
}

\begin{apient}
  EventStop::ptr getEventStop()
  EventStop::const_ptr getEventStop() const
\end{apient}

\apidesc{
}

\begin{apient}
  EventBreakpoint::ptr getEventBreakpoint()
  EventBreakpoint::const_ptr getEventBreakpoint() const
\end{apient}

\apidesc{
}

\begin{apient}
  EventNewThread::ptr getEventNewThread()
  EventNewThread::const_ptr getEventNewThread() const
\end{apient}

\apidesc{
}

\begin{apient}
  EventNewUserThread::ptr getEventNewUserThread()
  EventNewUserThread::const_tr getEventNewUserThread() const
\end{apient}

\apidesc{
}

\begin{apient}
  EventNewLWP::ptr getEventNewLWP()
  EventNewLWP::const_ptr getEventNewLWP() const
\end{apient}

\apidesc{
}

\begin{apient}
  EventThreadDestroy::ptr getEventThreadDestroy()
  EventThreadDestroy::const_ptr getEventThreadDestroy() const
\end{apient}

\apidesc{
}

\begin{apient}
  EventUserThreadDestroy::ptr getEventUserThreadDestroy()
  EventUserThreadDestroy::const_ptr getEventUserThreadDestroy() const
\end{apient}

\apidesc{
}

\begin{apient}
  EventLWPDestroy::ptr getEventLWPDestroy()
  EventLWPDestroy::const_ptr getEventLWPDestroy() const
\end{apient}

\apidesc{
}

\begin{apient}
  EventFork::ptr getEventFork()
  EventFork::const_ptr getEventFor() const
\end{apient}

\apidesc{
}

\begin{apient}
  EventSignal::ptr getEventSignal()
  EventSignal::const_ptr getEventSignal() const
\end{apient}

\apidesc{
}

\begin{apient}
  EventRPC::ptr getEventRPC()
  EventRPC::const_ptr getEventRPC() const
\end{apient}

\apidesc{
}

\begin{apient}
  EventSingleStep::ptr getEventSingleStep()
  EventSingleStep::const_ptr getEventSingleStep() const
\end{apient}

\apidesc{
}

\begin{apient}
  EventLibrary::ptr getEventLibrary()
  EventLibrary::const_ptr getEventLibrary() const
\end{apient}

\apidesc{
  These functions serve as a form of dynamic\_cast. They cast the Event into a
  child type and return the result of that cast. If the Event object is not of
  the appropriate type for the given function, then they return a shared pointer
  NULL equivalent (ptr() or const\_ptr()).
  For example, if an Event was an instance of an EventRPC, then the
  getEventRPC()
  function would cast it to EventRPC and return the resulting value.
}
